\section{The Probe Predicts Cohort EEG}
\label{sec:vaxis-brain}

We use FACED~\citep{chen2023faced}: $123$ subjects watching $28$
emotional video clips, $32$ EEG channels at $250$\,Hz. From the
cohort EEG averaged over subjects and time, we fit a single ridge
regression predicting the V-axis projection from the standard
$160$ channel--band features. The ridge prediction matches the
LLM-derived V-axis at $\mathbf{r{=}{+}0.87}$ ($n{=}28$ stimuli,
$p<10^{-9}$; Figure~\ref{fig:eeg-llm-circle})---one ridge model on
all features, not a multi-cell argmax. The headline correlation
holds against a battery of controls (matched-norm random-direction
null, split-half subject reliability, multi-comparison correction
over the channel--band grid, and class-level ranking) catalogued in
Appendix~\ref{app:brain-deep-dive}.

\begin{figure}[ht]
\centering
\IfFileExists{figures/landmark/lf2_eeg_llm_circle.pdf}
  {\includegraphics[width=\linewidth]{landmark/lf2_eeg_llm_circle.pdf}}
  {\includegraphics[width=\linewidth]{F2_eeg_llm_circle.pdf}}
\caption[EEG-LLM cohort circle]{\textbf{Category-level alignment between the LLM-derived valence
direction and human EEG.}
Each point is one of the $28$ FACED video stimuli; $x$ is the
\emph{CLIP-text} V-axis value of each stimulus's emotion category (one
value per class),
$y$ is the single ridge-regression prediction of that value from
all $160$ cohort channel--band features ($32$ channels $\times$ $5$
frequency bands, averaged over $123$ subjects and each clip).
Cohort Pearson $r{=}{+}0.87$ at emotion-category resolution
($n{=}28$, nested-perm $p<2{\times}10^{-4}$; per-clip $r{=}0.71$); the nine-story
Qwen3-4B V-axis reaches $r{=}0.80$ on the same ridge. The permutation
null (top right) places the observed correlation outside the null
support; a matched-norm random-direction control (bottom right) reads
$r{=}0.07$. The axis is extracted with no EEG supervision.
}
\label{fig:eeg-llm-circle}
\end{figure}

The same picture survives across all $14$ LLMs: the family-tier
structure recapitulates on the brain side, with the Qwen3 family
alone reaching mean $r{=}{+}0.796 \pm 0.011$ irrespective of size.
The FACED-derived V-axis also generalises to a held-out EEG dataset
without retraining, ranking the five SEED-V emotion
classes~\citep{liu2022seedv} at semantic-space $r{=}{+}0.96$;
Appendix~\ref{app:seedv-full} re-derives the V-axis from scratch on
SEED-V and re-tests every claim in the paper.
